\newpage

\section{量子力学：平面转子在电场中的简并微扰论}
\label{sec:planar-rotator}

\textbf{题目：平面转子在匀强电场中的微扰修正}

质量为 $m$、带电量为 $-e$ 的粒子被限制在 $X$-$Y$ 平面内，绕原点做半径为 $r$ 的圆周运动。求该粒子做圆周运动的能量本征值和本征函数。

若在平面内沿 $-x$ 方向施加匀强电场 $\vec{\mathscr{E}} = -\mathscr{E}\,\vec{e}_x$（$\mathscr{E}>0$），运用微扰理论，求第一激发态与第二激发态能量的二级近似。

\subsection{考点快速分析}

\begin{itemize}
    \item \textbf{角动量代数}：$\hat{L}_z = -i\hbar\frac{\partial}{\partial\phi}$，周期性边界条件 $\psi(\phi+2\pi)=\psi(\phi)$
    \item \textbf{简并结构}：基态 $m=0$ 非简并；激发态 $|m|\geq 1$ 二重简并（$m$ 与 $-m$ 能量相同）
    \item \textbf{微扰矩阵元}：$H' = er\mathscr{E}\cos\phi$，选择定则 $\Delta m = \pm 1$
    \item \textbf{简并微扰的陷阱}：简并子空间内矩阵元全为零，一级修正消失
    \item \textbf{非简并公式的合法性}：$H'$ 不耦合简并子空间内的态，可直接套用二级微扰公式
\end{itemize}

\subsection{未微扰系统的本征值与本征函数}

转动惯量 $I = mr^2$，哈密顿量为
\[
\hat{H}_0 = \frac{\hat{L}_z^2}{2I} = -\frac{\hbar^2}{2mr^2}\frac{\partial^2}{\partial\phi^2}.
\]
本征方程 $\hat{H}_0\psi_m = E_m^{(0)}\psi_m$ 的解满足周期性边界条件 $\psi(\phi+2\pi)=\psi(\phi)$：
\[
\boxed{\,\psi_m(\phi) = \frac{1}{\sqrt{2\pi}}e^{im\phi},\quad m = 0, \pm 1, \pm 2, \dots\,}
\]
\[
\boxed{\,E_m^{(0)} = \frac{m^2\hbar^2}{2mr^2} = \frac{m^2\hbar^2}{2I}\,}.
\]

\textbf{简并结构}：
\begin{itemize}
    \item $m=0$（基态）：$E_0^{(0)}=0$，\textbf{非简并}
    \item $|m|\geq 1$（激发态）：$E_m^{(0)} = E_{-m}^{(0)}$，\textbf{二重简并}
\end{itemize}

\subsection{微扰项的矩阵元与选择定则}

电场沿 $-x$ 方向，微扰势能：
\[
H' = -q\vec{\mathscr{E}}\cdot\vec{r} = (-e)(-\mathscr{E})(r\cos\phi) = er\mathscr{E}\cos\phi.
\]
计算矩阵元（利用 $\cos\phi = \frac{e^{i\phi}+e^{-i\phi}}{2}$）：
\begin{align*}
\langle m'|H'|m\rangle &= er\mathscr{E}\int_0^{2\pi}\frac{e^{-im'\phi}}{\sqrt{2\pi}}\cos\phi\,\frac{e^{im\phi}}{\sqrt{2\pi}}\,d\phi\\
&= \frac{er\mathscr{E}}{4\pi}\int_0^{2\pi}\bigl[e^{i(m-m'+1)\phi} + e^{i(m-m'-1)\phi}\bigr]d\phi\\
&= \frac{er\mathscr{E}}{2}\bigl(\delta_{m',m+1} + \delta_{m',m-1}\bigr).
\end{align*}

\textbf{选择定则：为什么只有 $\Delta m = \pm 1$？}

微扰 $H' \propto \cos\phi$ 是偶函数，在角动量表象中相当于一个\textbf{升降算符的组合}：$\cos\phi \sim e^{i\phi} + e^{-i\phi}$，每次作用使磁量子数改变 $\pm 1$。这类似于谐振子问题中 $x \sim a + a^\dagger$ 只耦合相邻能级。

\textbf{关键推论}：对简并子空间 $\{|m\rangle, |-m\rangle\}$（$m\geq 1$），
\[
\langle m|H'|-m\rangle = \frac{er\mathscr{E}}{2}(\delta_{m,-m+1} + \delta_{m,-m-1}) = 0
\]
因为 $2m = 1$ 或 $2m = -1$ 对整数 $m$ 均不成立。所以\textbf{简并子空间内的微扰矩阵元全为零}。

\subsection{定态微扰理论：二阶修正——确定公式与记忆锚点}

本题的难点不在于代数运算，而在于\textbf{二级修正的公式本身}。考场高压下极易记混分母顺序、漏掉模方、或混淆简并/非简并的适用条件。下面给出最常用、最确定的公式，并标注每一个系数的来源。

\textbf{非简并微扰：二阶能量修正（最核心公式）}

设 $H = H_0 + H'$，$H_0|n^{(0)}\rangle = E_n^{(0)}|n^{(0)}\rangle$，且能级 $E_n^{(0)}$ \textbf{非简并}。则能量本征值到二级近似为
\[
E_n = E_n^{(0)} + E_n^{(1)} + E_n^{(2)} + O\bigl((H')^3\bigr),
\]
其中各级修正为
\begin{align*}
E_n^{(1)} &= \langle n^{(0)}|H'|n^{(0)}\rangle, \\
\boxed{\,E_n^{(2)} = \sum_{m\neq n} \frac{|\langle m^{(0)}|H'|n^{(0)}\rangle|^2}{E_n^{(0)} - E_m^{(0)}}\,}&.
\end{align*}

\textbf{符号与单位制约定}（无一例外，所有教材通用）：
\begin{itemize}
    \item 分子是 $\displaystyle |H'_{mn}|^2 = |\langle m|H'|n\rangle|^2$，\textbf{必须取模方}，保证结果恒为实数；
    \item 分母是 $\displaystyle E_n^{(0)} - E_m^{(0)}$，\textbf{待修正能级在分子位置、在中间态前面}——记作 ``$n$ 在前''；
    \item 求和跑遍\textbf{除 $n$ 以外的所有中间态}（包括简并子空间外的态，以及简并子空间内与 $n$ 不同的态）；
    \item 没有额外的 $1/2$、$\pm i$ 等系数，公式\textbf{exactly 如上}。
\end{itemize}

\textbf{分母顺序是第一大杀手——三种锚定方法}

\textbf{核心问题}：容易把分母记成 $E_m^{(0)} - E_n^{(0)}$（中间态减待求态）。下面给出三种互不依赖的记忆锚，任选其一即可：

\textbf{锚定法 1：``求谁谁在前''}

你要求的是 $E_n^{(2)}$（第 $n$ 个能级的修正），分母就必须是 $E_n^{(0)} - E_m^{(0)}$——\textbf{待求态的能量在前，中间态的能量在后}。和一级波函数修正的分母完全一致：
\[
|n^{(1)}\rangle = \sum_{m\neq n}\frac{\langle m|H'|n\rangle}{E_n^{(0)}-E_m^{(0)}}|m\rangle,
\quad
E_n^{(2)} = \langle n|H'|n^{(1)}\rangle.
\]
二级能量公式本质上就是 ``$\langle n|H'$ 作用在一级波函数上''，分母自然继承一级波函数的分母。

\textbf{锚定法 2：基态符号检验（物理直觉）}

若 $H'$ 使系统\textbf{基态能量降低}（如感生电偶极矩与外场反向，体系被极化后能量下降），则基态的二级修正必须是\textbf{负数}。对基态 $n=0$，$E_0^{(0)} < E_m^{(0)}$，分母 $E_0^{(0)}-E_m^{(0)}<0$，分子 $|H'_{m0}|^2>0$，故 $E_0^{(2)}<0$，符合直觉。若分母记反了，$E_0^{(2)}>0$，基态能量反而升高——\textbf{立即发现错误}。

\textbf{锚定法 3：分母 = ``初态能量 $-$ 虚中间态能量''}

把二级修正想象成 ``从 $n$ 出发，虚跃迁到 $m$，再返回 $n$''。分母反映的是\textbf{虚中间态与初态的能量失配}。初态是 $n$，所以 $E_n^{(0)}$ 必须在前面。

\textbf{公式从哪里来？——二级微扰的骨架推导（考场可省略）}

将 $|n\rangle = |n^{(0)}\rangle + |n^{(1)}\rangle + \cdots$ 和 $E_n = E_n^{(0)} + E_n^{(1)} + E_n^{(2)} + \cdots$ 代入定态方程 $(H_0+H')|n\rangle = E_n|n\rangle$，按 $(H')$ 的幂次收集：
\begin{itemize}
    \item $O(H'^0)$：$H_0|n^{(0)}\rangle = E_n^{(0)}|n^{(0)}\rangle$（已知）。
    \item $O(H'^1)$：$H_0|n^{(1)}\rangle + H'|n^{(0)}\rangle = E_n^{(0)}|n^{(1)}\rangle + E_n^{(1)}|n^{(0)}\rangle$。
    \item $O(H'^2)$：$H_0|n^{(2)}\rangle + H'|n^{(1)}\rangle = E_n^{(0)}|n^{(2)}\rangle + E_n^{(1)}|n^{(1)}\rangle + E_n^{(2)}|n^{(0)}\rangle$。
\end{itemize}
对 $O(H'^2)$ 方程左乘 $\langle n^{(0)}|$，利用 $\langle n^{(0)}|H_0 = E_n^{(0)}\langle n^{(0)}|$ 和 $E_n^{(1)} = \langle n^{(0)}|H'|n^{(0)}\rangle$，得
\[
E_n^{(2)} = \langle n^{(0)}|H'|n^{(1)}\rangle - E_n^{(1)}\langle n^{(0)}|n^{(1)}\rangle.
\]
一级修正的波函数为 $|n^{(1)}\rangle = \sum_{m\neq n} \frac{\langle m|H'|n\rangle}{E_n^{(0)}-E_m^{(0)}}|m^{(0)}\rangle$，且通常取 $\langle n^{(0)}|n^{(1)}\rangle = 0$（相位约定）。代入即得上述 $E_n^{(2)}$ 公式。

\textbf{简并微扰中何时能直接套用非简并二级公式？}

对简并子空间 $\mathcal{D}_n = \{|n,\alpha\rangle\}_{\alpha=1}^{g_n}$（$g_n$ 重简并），先求一级微扰矩阵
\[
W_{\alpha\beta} = \langle n,\alpha|H'|n,\beta\rangle.
\]
\textbf{若 $W$ 在简并子空间内恒为零}（如本题 $\langle m|H'|-m\rangle=0$），则：
\begin{enumerate}
    \item 一级修正 $E_n^{(1)}=0$；
    \item 二级修正\textbf{可直接使用非简并公式}，只需将求和中的 ``$m\neq n$'' 理解为 ``$m$ 跑遍\textbf{所有}不在 $|n\rangle$ 上的态''（包括简并子空间内的伴生态和子空间外的所有态）。
\end{enumerate}
\textbf{物理原因}：二级修正的本质是 ``$n$ 先通过 $H'$ 跳到某个中间态 $m$，再从 $m$ 跳回 $n$''。只要中间过程不依赖于简并子空间内部如何选取基矢，结果就与基的选取无关，非简并公式自然成立。

\textbf{波函数的二阶修正（完整版，考场视要求选用）}

若题目要求波函数也修正到二级，公式为
\[
|n^{(2)}\rangle = \sum_{m\neq n}|m^{(0)}\rangle\Biggl[
\sum_{\substack{l\neq n}} \frac{\langle m|H'|l\rangle\langle l|H'|n\rangle}{(E_n^{(0)}-E_m^{(0)})(E_n^{(0)}-E_l^{(0)})}
- \frac{\langle n|H'|n\rangle\langle m|H'|n\rangle}{(E_n^{(0)}-E_m^{(0)})^2}
- \frac{\langle m|H'|n\rangle\langle n|H'|n\rangle}{(E_n^{(0)}-E_m^{(0)})^2}
\Biggr].
\]
在 $\langle n^{(0)}|n^{(2)}\rangle=0$ 的相位约定下可化简为常见教材形式。\textbf{对能量计算而言，只需记住 $E_n^{(2)}$ 的公式即可}；波函数二级修正只在计算跃迁矩阵元或\textit{expectation values}时才需要。

\textbf{考场速记：二阶微扰的「三字诀」}
\begin{enumerate}
    \item \textbf{``方''}——分子必须是 $|H'_{mn}|^2$（模方，恒正）；
    \item \textbf{``前''}——分母是 $E_n^{(0)} - E_m^{(0)}$（待求态 $n$ 在前）；
    \item \textbf{``跑''}——求和跑遍所有 $m\neq n$（不遗漏中间态）。
\end{enumerate}
把三字诀念三遍：\textbf{方、前、跑}。动笔前先在草稿纸右上角写下
\[
E_n^{(2)} = \sum_{m\neq n}\frac{|\langle m|H'|n\rangle|^2}{E_n^{(0)}-E_m^{(0)}}
\]
作为「护身符」，后续所有数值代入都围绕它展开，可大幅降低记忆失误。

\subsection{简并微扰的处理}

对 $|m|\geq 1$ 的简并子空间，微扰矩阵为
\[
W = \begin{pmatrix}
\langle m|H'|m\rangle & \langle m|H'|-m\rangle\\[4pt]
\langle -m|H'|m\rangle & \langle -m|H'|-m\rangle
\end{pmatrix} = \begin{pmatrix} 0 & 0 \\ 0 & 0 \end{pmatrix}.
\]

因此：
\begin{enumerate}
    \item \textbf{一级修正}：$E_m^{(1)} = 0$（简并微扰的一级修正为零）。
    \item \textbf{二级修正}：由于 $H'$ 不耦合简并子空间内的态，$|m\rangle$ 和 $|-m\rangle$ 在二级微扰中各自独立演化，\textbf{可以直接套用非简并微扰的二级公式}：
    \[
    E_m^{(2)} = \sum_{k\neq m}\frac{|\langle k|H'|m\rangle|^2}{E_m^{(0)} - E_k^{(0)}}.
    \]
\end{enumerate}

\subsection{第一激发态 $|m|=1$ 的二级修正}

对 $|1\rangle$，非零矩阵元连接到 $|0\rangle$ 和 $|2\rangle$：
\[
E_1^{(2)} = \frac{|\langle 0|H'|1\rangle|^2}{E_1^{(0)}-E_0^{(0)}} + \frac{|\langle 2|H'|1\rangle|^2}{E_1^{(0)}-E_2^{(0)}}.
\]
代入 $|\langle 0|H'|1\rangle| = |\langle 2|H'|1\rangle| = \frac{er\mathscr{E}}{2}$ 和能量值：
\begin{align*}
E_1^{(2)} &= \frac{(er\mathscr{E}/2)^2}{\frac{\hbar^2}{2mr^2}-0} + \frac{(er\mathscr{E}/2)^2}{\frac{\hbar^2}{2mr^2}-\frac{4\hbar^2}{2mr^2}}\\
&= \frac{m r^4 e^2 \mathscr{E}^2}{2\hbar^2} - \frac{m r^4 e^2 \mathscr{E}^2}{6\hbar^2}\\
&= \boxed{\,\frac{m r^4 e^2 \mathscr{E}^2}{3\hbar^2}\,}.
\end{align*}

由对称性，$|-1\rangle$ 给出完全相同的二级修正。

\subsection{第二激发态 $|m|=2$ 的二级修正}

对 $|2\rangle$，非零矩阵元连接到 $|1\rangle$ 和 $|3\rangle$：
\[
E_2^{(2)} = \frac{|\langle 1|H'|2\rangle|^2}{E_2^{(0)}-E_1^{(0)}} + \frac{|\langle 3|H'|2\rangle|^2}{E_2^{(0)}-E_3^{(0)}}.
\]
代入：
\begin{align*}
E_2^{(2)} &= \frac{(er\mathscr{E}/2)^2}{\frac{4\hbar^2}{2mr^2}-\frac{\hbar^2}{2mr^2}} + \frac{(er\mathscr{E}/2)^2}{\frac{4\hbar^2}{2mr^2}-\frac{9\hbar^2}{2mr^2}}\\
&= \frac{m r^4 e^2 \mathscr{E}^2}{6\hbar^2} - \frac{m r^4 e^2 \mathscr{E}^2}{10\hbar^2}\\
&= \boxed{\,\frac{m r^4 e^2 \mathscr{E}^2}{15\hbar^2}\,}.
\end{align*}

\begin{center}
\renewcommand{\arraystretch}{1.4}
\begin{tabular}{|c|c|c|c|c|}
\hline
\textbf{能级} & \textbf{状态} & \textbf{简并度} & \boldsymbol{$E^{(1)}$} & \boldsymbol{$E^{(2)}$} \\
\hline
基态 ($m=0$) & $|0\rangle$ & 1 & $0$ & $-\dfrac{mr^4e^2\mathscr{E}^2}{\hbar^2}$ \\
\hline
第一激发态 ($|m|=1$) & $|{\pm}1\rangle$ & 2 & $0$ & $\dfrac{mr^4e^2\mathscr{E}^2}{3\hbar^2}$ \\
\hline
第二激发态 ($|m|=2$) & $|{\pm}2\rangle$ & 2 & $0$ & $\dfrac{mr^4e^2\mathscr{E}^2}{15\hbar^2}$ \\
\hline
\end{tabular}
\end{center}

\subsection{基态的二级修正（补充）}

虽然题目只要求激发态，但基态 ($m=0$) 的修正也很有教学意义。$|0\rangle$ 只与 $|{\pm}1\rangle$ 耦合：
\[
E_0^{(2)} = \frac{|\langle 1|H'|0\rangle|^2}{E_0^{(0)}-E_1^{(0)}} + \frac{|\langle -1|H'|0\rangle|^2}{E_0^{(0)}-E_{-1}^{(0)}}
= 2\times\frac{(er\mathscr{E}/2)^2}{0-\frac{\hbar^2}{2mr^2}}
= -\frac{mr^4e^2\mathscr{E}^2}{\hbar^2}.
\]
负号符合物理直觉：感生电偶极矩与外电场反向，能量降低。
